\section{The foundational debate: three classical positions}
\label{sec:debate}

Since the founding paper of Aharonov and Bohm
\cite{AharonovBohm1959}, quantum electrodynamics has posed a
foundational question that resists any clean experimental
resolution: which among the 4-potential $A_\mu$, the field tensor
$F_{\mu\nu}$, and the holonomy phase $\oint_\gamma A \cdot d\ell$
should be regarded as the physically real object? Three
positions divide the contemporary literature, with no classical
experiment so far able to distinguish them. This section recalls
them in their sharpest form, for the sole purpose of fixing what
Persistence Theory (PT) claims to adjudicate. The review is not
exhaustive --- for the philosophical developments we refer to the
monographs of Healey \cite{Healey2007} and Maudlin
\cite{Maudlin2011}.

\subsection{Position A: reality of the vector potential}
\label{ssec:position_A_en}

Position A, inherited from the Bohm--Hiley programme
\cite{BohmHiley1993}, holds that the potential $A_\mu$ is
physically real and that the tensor $F_{\mu\nu} = \partial_\mu A_\nu
- \partial_\nu A_\mu$ is a derived quantity. The empirical argument
is the Aharonov--Bohm effect~\cite{AharonovBohm1959}: the
interferometric phase shift $\Delta\varphi = (e/\hbar) \oint A \cdot
d\ell$ recorded by a particle that loops around an ideal solenoid
is measurable even where $E = B = 0$ along its entire trajectory.
No field crosses the path; no local object can therefore account
for the effect. The potential, by contrast, is non-zero (up to
gauge) along the trajectory, and the Bohmian programme reads this
as proof that $A_\mu$ must be regarded as an object of higher
ontological status.

The known difficulty of position A is gauge dependence. The
potential $A_\mu$ is defined only up to a transformation
$A_\mu \mapsto A_\mu + \partial_\mu \Lambda$, and every physically
measurable quantity must be invariant under this transformation.
The Aharonov--Bohm phase is gauge-invariant \emph{only as a closed
integral}; its local values are not. Position A must therefore
either accept a physical reality that is not itself gauge-invariant
--- which many consider a major ontological concession --- or
treat $A_\mu$ as a hybrid object only the gauge-equivalence class
of which is real.

\subsection{Position B: intrinsic non-locality of the field}
\label{ssec:position_B_en}

Position B accepts $F_{\mu\nu}$ as the real object but makes it
intrinsically non-local. In the strong reading proposed by Mead
\cite{Mead2000} and in the wake of the analyses of Bell
\cite{Bell1964} and the experimental tests of Aspect
\cite{Aspect1982} and Hensen et al.~\cite{Hensen2015},
the electromagnetic field is not a distribution on classical
spacetime: it carries a correlation structure between spatially
separated points that is not explained by the local variation of
its components. The Aharonov--Bohm phase then becomes the signal
of such non-locality, and the EPR correlations its extension to
matter.

The structural objection to position B is that it requires
modifying the classical ontology of the field to make it
intrinsically non-local, without specifying which mathematical
mechanism implements the modification. The available formal
machinery (connection fibre bundles, gauge theories, Wilson
operators) admits strictly local interpretations under a suitable
choice of variables. The non-locality of position B remains, in
most formulations, a property of the \emph{interpretation} of the
field rather than of its mathematical structure.

\subsection{Position C: primacy of holonomy}
\label{ssec:position_C_en}

Position C, due to Wu and Yang \cite{WuYang1975} and developed by
Healey \cite{Healey2007} and Belot \cite{Belot1998}, holds that
neither of the two local objects $A_\mu, F_{\mu\nu}$ is
fundamental: $A_\mu$ because it depends on gauge, $F_{\mu\nu}$
because it does not contain the Aharonov--Bohm phase information
in regions where it vanishes. The only object that is at once
gauge-invariant and sensitive to the Aharonov--Bohm experiment is
the Wilson loop $W_\gamma = \exp\bigl[i (e/\hbar) \oint_\gamma A
\cdot d\ell\bigr]$, a non-local functional of closed paths $\gamma$.
This object, and not $A_\mu$ or $F_{\mu\nu}$, deserves the status
of primitive.

Position C is conceptually satisfying --- it dissolves the
asymmetry between the gauge-dependence of $A_\mu$ and the
incompleteness of $F_{\mu\nu}$ --- but it leaves open the question
of the \emph{origin} of Wilson loops: why these particular
non-local objects and not others? Classical gauge theory introduces
them as a post hoc mathematical construction from a connection
that remains itself primary in the formalism. Position C therefore
remains, in a sense, a description rather than an explanation.

\subsection{The operational impasse}
\label{ssec:impasse_en}

Positions A, B, C make the same experimental predictions within
standard quantum electrodynamics. The Aharonov--Bohm phase is
measurable in all three frameworks; the EPR/Bell correlations are
reproduced in all three; cross sections and lifetimes are
identical. No measurement to date has been able to distinguish
the three positions at the level of quantitative prediction. The
debate is thus, at present, a debate of interpretation.

Persistence Theory modifies this state of affairs on two counts.
On one hand, it promotes position C to a \emph{theorem} rather
than an interpretation: the modular holonomy angle $\theta_p$ is
the unique primitive invariant satisfying four structural
conditions that neither $A_\mu$ nor $F_{\mu\nu}$ can meet
simultaneously (Section~\ref{sec:holonomy_primitive}). On the
other hand, it produces an experimental prediction --- prediction
QH1a of Section~\ref{sec:qh1a_prediction} --- whose empirical
verdict effectively distinguishes PT (and with it position C)
from a standard theory on curved manifold. The foundational
debate is no longer merely a debate of interpretation: under PT,
it becomes adjudicable.
